\section{2025年高考题}

\subsection{2025年全国I卷}\label{gk:2025年全国I卷}

\begin{description}
    \item[一、选择题] 本题共8小题，每小题5分，共40分．在每小题给出的四个选项中，只有一项是符合题目要求的．
    \begin{enumerate}[leftmargin=5pt]
        \item $(1+5\mathrm{i})\mathrm{i}$ 的虚部为
        
        {\centering\begin{tabular*}{\linewidth}{@{}@{\extracolsep{\fill}}llll@{}}
          \(\text{A. }-1\) & \(\text{B. }0\) & \(\text{C. }1\) & \(\text{D. }6\)
        \end{tabular*}\par}

        \item 设全集 $U=\{x\mid x\text{为小于}9\text{的正整数}\}$，集合 $A=\{1,3,5\}$，则 $\complement_U A$ 中元素的个数为
        
        {\centering\begin{tabular*}{\linewidth}{@{}@{\extracolsep{\fill}}llll@{}}
          \(\text{A. }0\) & \(\text{B. }3\) & \(\text{C. }5\) & \(\text{D. }8\)
        \end{tabular*}\par}

        \item 双曲线 $C$ 的虚轴长为实轴长的 $7$ 倍，则 $C$ 的离心率为
       
        {\centering\begin{tabular*}{\linewidth}{@{}@{\extracolsep{\fill}}llll@{}}
          \(\text{A. }2\) & \(\text{B. }\sqrt{2}\) & \(\text{C. }7\) & \(\text{D. }2\sqrt{2}\)
        \end{tabular*}\par}

        \item 已知点 $(a,0)\,(a>0)$ 是函数 $y=2\tan\bigl(x-\frac{\pi}{3}\bigr)$ 的图像的一个对称中心，则 $a$ 的最小值为
       
        {\centering\begin{tabular*}{\linewidth}{@{}@{\extracolsep{\fill}}llll@{}}
          \(\text{A. }\dfrac{\pi}{6}\) & \(\text{B. }\dfrac{\pi}{3}\) & \(\text{C. }\dfrac{\pi}{2}\) & \(\text{D. }\dfrac{2\pi}{3}\)
        \end{tabular*}\par}

        \item 已知 $f(x)$ 是定义在 $\mathbb{R}$ 上且周期为 $2$ 的偶函数，当 $2\leqslant x\leqslant3$ 时，$f(x)=5-2x$，则 $f\bigl(-\frac{3}{4}\bigr)=$
     
        {\centering\begin{tabular*}{\linewidth}{@{}@{\extracolsep{\fill}}llll@{}}
          \(\text{A. }-\dfrac{1}{2}\) & \(\text{B. }-\dfrac{1}{4}\) & \(\text{C. }\dfrac{1}{4}\) & \(\text{D. }\dfrac{1}{2}\)
        \end{tabular*}\par}

        \item 帆船比赛中，运动员可借助风力计测定风速的大小和方向，测出的结果在航海学中称为视风风速，视风风速对应的向量，是真风风速对应的向量与船行风速对应的向量之和，其中船行风速对应的向量与船速对应的向量大小相等，方向相反．图1给出了部分风力等级、名称与风速大小的对应关系．已知某帆船运动员在某时刻测得的视风风速对应的向量与船速对应的向量如图2（风速的大小和向量的大小相同，单位：$\text{m/s}$），则真风为
      
        {\centering\begin{tabular*}{\linewidth}{@{}@{\extracolsep{\fill}}llll@{}}
          \(\text{A. }\)轻风 & \(\text{B. }\)微风 & \(\text{C. }\)和风 & \(\text{D. }\)劲风
        \end{tabular*}\par}

        \item 已知圆 $x^2+(y+2)^2=r^2\,(r>0)$ 上到直线 $y=3x+2$ 的距离为 $1$ 的点有且仅有 $2$ 个，则 $r$ 的取值范围是
       
        {\centering\begin{tabular*}{\linewidth}{@{}@{\extracolsep{\fill}}llll@{}}
          \(\text{A. }(0,1)\) & \(\text{B. }(1,3)\) & \(\text{C. }(3,+\infty)\) & \(\text{D. }(0,+\infty)\)
        \end{tabular*}\par}

        \item 若 $2+\log_2 x=3+\log_3 y=5+\log_5 z$，则 $x,y,z$ 的大小关系不可能是
     
        {\centering\begin{tabular*}{\linewidth}{@{}@{\extracolsep{\fill}}llll@{}}
          \(\text{A. }x>y>z\) & \(\text{B. }x>z>y\) & \(\text{C. }y>x>z\) & \(\text{D. }y>z>x\)
        \end{tabular*}\par}
    \end{enumerate}
\end{description}

\begin{description}
    \item[二、选择题] 本题共3小题，每小题6分，共18分．在每小题给出的选项中，有多项符合题目要求．全部选对的得6分，部分选对的得部分分，有选错的得0分．
    \begin{enumerate}[leftmargin=5pt]
      \addtocounter{enumi}{+8}
        \item 在正三棱柱 $ABC-A_1B_1C_1$ 中，$D$ 为 $BC$ 中点，则
        \begin{enumerate}[label=\Alph*.]
            \item $AD\perp A_1C$
            \item $BC\perp$ 平面 $AA_1D$
            \item $AD\parallel A_1B_1$
            \item $CC_1\parallel$ 平面 $AA_1D$
        \end{enumerate}

        \item 设抛物线 $C:y^2=6x$ 的焦点为 $F$，过 $F$ 的直线与 $C$ 交于 $A,B$ 两点，过 $A$ 作直线 $l:x=-\frac{3}{2}$ 的垂线，垂足为 $D$，过 $F$ 且垂直于 $AB$ 的直线与 $l$ 交于点 $E$，则
        \begin{enumerate}[label=\Alph*.]
            \item $|AD|=|AF|$
            \item $|AB|=|AE|$
            \item $|AB|\geqslant6$
            \item $|AE|\cdot|BE|\geqslant18$
        \end{enumerate}

        \item 已知 $\triangle ABC$ 的面积为 $\dfrac{1}{4}$，若 $\cos2A+\cos2B+2\sin C=2$，$\cos A\cos B\sin C=\dfrac{1}{4}$，则
        \begin{enumerate}[label=\Alph*.]
            \item $\sin C=\sin2A+\sin2B$
            \item $AB=\sqrt{2}$
            \item $\sin A+\sin B=\dfrac{\sqrt{6}}{2}$
            \item $AC^2+BC^2=3$
        \end{enumerate}
    \end{enumerate}
\end{description}

\begin{description}
    \item[三、填空题] 本题共3小题，每小题5分，共15分．
    \begin{enumerate}[leftmargin=5pt]
      \addtocounter{enumi}{+11}
        \item 若直线 $y=2x+5$ 是曲线 $y=\mathrm{e}^x+x+a$ 的切线，则 $a=$ \rule[-0.3ex]{3em}{0.4pt}．
        \item 若一个等比数列的各项均为正数，前 $4$ 项和等于 $4$，前 $8$ 项和等于 $68$，则该等比数列的公比等于 \rule[-0.3ex]{3em}{0.4pt}．
        \item 一个箱子里有 $5$ 个相同的球，分别标有数字\(1,2,3,4,5\)，
        从中有放回地取\(3\)次，每次取\(1\)个球. 记 $X$为这\(5\)个球中至少被取出\(1\)次的球的个数，则 $\mathrm{E}(X)=$ \rule[-0.3ex]{3em}{0.4pt}．
    \end{enumerate}
\end{description}

\begin{description}
    \item[四、解答题] 本题共5小题，共77分．解答应写出文字说明、证明过程或演算步骤．
    \begin{enumerate}[leftmargin=5pt]
      \addtocounter{enumi}{+14}
        \item （13分）

        为研究某疾病与超声波检查结果的关系，从做过超声波检查的人群中随机调查了 $1\,000$ 人，得到如下列联表．
        \begin{table}[H]
        \centering
        \begin{tabular}{|c|c|c|c|}
        \hline
        \textbf{组别} & \textbf{正常} & \textbf{不正常} & \textbf{合计} \\
        \hline
        患该疾病 & 20 & 180 & 200 \\
        \hline
        未患该疾病 & 780 & 20 & 800 \\
        \hline
        合计 & 800 & 200 & 1\,000 \\
        \hline
        \end{tabular}
        \end{table}

        \begin{enumerate}[label=(\arabic*)]
            \item 记超声波检查结果不正常者患该疾病的概率为 $p$，求 $p$ 的估计值；
            \item 依据小概率值 $\alpha=0.001$ 的独立性检验，分析超声波检查结果是否与患该疾病有关．
            \begin{quote}
            附：$K^2=\dfrac{n(ad-bc)^2}{(a+b)(c+d)(a+c)(b+d)}$，\quad
            $\begin{array}{c|ccc}
            P(K^2\geqslant k) & 0.050 & 0.010 & 0.001 \\
            \hline
            k & 3.841 & 6.635 & 10.828
            \end{array}$．
            \end{quote}
        \end{enumerate}

        \item （15分）

        已知数列 $\{a_n\}$ 满足 $a_1=3$，$\dfrac{a_{n+1}}{n}=\dfrac{a_n}{n+1}+\dfrac{1}{n(n+1)}$．

        \begin{enumerate}[label=(\arabic*)]
            \item 证明：数列 $\{na_n\}$ 是等差数列；
            \item 给定正整数 $m$，设函数 $f(x)=a_1x+a_2x^2+\cdots+a_mx^m$，求 $f'(-2)$．
        \end{enumerate}

        \item （15分）

        如图，四棱锥 $P-ABCD$ 中，$PA\perp$ 底面 $ABCD$，$BC\parallel AD$，$AB\perp AD$．

        \begin{enumerate}[label=(\arabic*)]
            \item 证明：平面 $PAB\perp$ 平面 $PAD$；
            \item 若 $PA=AB=\sqrt{2}$，$AD=1+\sqrt{3}$，$BC=2$，$P,B,C,D$ 四点在同一个球面上，设该球面的球心为 $O$．
            \begin{enumerate}[label=(\roman*)]
                \item 证明：$O$ 在平面 $ABCD$ 上；
                \item 求直线 $AC$ 与直线 $PO$ 所成角的余弦值．
            \end{enumerate}
        \end{enumerate}

        \item （17分）

        已知椭圆 $C:\dfrac{x^2}{a^2}+\dfrac{y^2}{b^2}=1\,(a>b>0)$ 的离心率为 $\dfrac{2\sqrt{2}}{3}$，下顶点为 $A$，右顶点为 $B$，且 $|AB|=\sqrt{10}$．

        \begin{enumerate}[label=(\arabic*)]
            \item 求 $C$ 的方程；
            \item 已知动点\(P\)不在 $y$ 轴上，点$R$ 在射线 $AP$ 上，且满足\(|AP|\cdot|AR|=3\)．
            \begin{enumerate}[label=(\roman*)]
                \item 设$P(m,n)$ ，求 $R$ 的坐标；（用 $m,n$ 表示）
                \item $O$ 为坐标原点，$Q$ 是 $C$ 上的动点，若直线 $OR$ 的斜率是直线 $OP$ 的斜率的 $3$ 倍，求 $|PQ|$ 的最大值．
            \end{enumerate}
        \end{enumerate}

        \item （17分）

        \begin{enumerate}[label=(\arabic*)]
            \item 求函数 $f(x)=5\cos x-\cos5x$ 在区间 $\bigl[0,\frac{\pi}{4}\bigr]$ 的最大值；
            \item 给定 $\theta\in(0,\pi)$ 和 $a\in\mathbb{R}$，证明：存在 $y\in[a-\theta,a+\theta]$，使得 $\cos y\leqslant\cos\theta$；
            \item 设\(b\in\mathbb{R}\), 若存在 $\varphi$，使得 $5\cos x-\cos(5x+\varphi)\leqslant b$对任意 $x\in \mathbb{R}$ 恒成立，求 $b$ 的最小值．
        \end{enumerate}
    \end{enumerate}
\end{description}

\clearpage


\subsection{2025年全国II卷}\label{gk:2025年全国II卷}

\begin{description}
    \item[一、单选题] 本题共8小题，每小题5分，共40分．每小题只有一个选项符合要求．
    \begin{enumerate}[leftmargin=5pt]
        \item 样本数据 $2,8,14,16,20$ 的平均数为
       
        {\centering\begin{tabular*}{\linewidth}{@{}@{\extracolsep{\fill}}llll@{}}
          \(\text{A. }8\) & \(\text{B. }9\) & \(\text{C. }12\) & \(\text{D. }18\)
        \end{tabular*}\par}

        \item 已知 $z=1+\mathrm{i}$，则 $\dfrac{1}{z-1}=$
      
        {\centering\begin{tabular*}{\linewidth}{@{}@{\extracolsep{\fill}}llll@{}}
          \(\text{A. }-\mathrm{i}\) & \(\text{B. }\mathrm{i}\) & \(\text{C. }-1\) & \(\text{D. }1\)
        \end{tabular*}\par}

        \item 已知集合 $A=\{-4,0,1,2,8\}$，$B=\{x\mid x^3=x\}$，则 $A\cap B=$
      
        {\centering\begin{tabular*}{\linewidth}{@{}@{\extracolsep{\fill}}llll@{}}
          \(\text{A. }\{0,1,2\}\) & \(\text{B. }\{1,2,8\}\) & \(\text{C. }\{2,8\}\) & \(\text{D. }\{0,1\}\)
        \end{tabular*}\par}

        \item 不等式 $\dfrac{x-4}{x-1}\geqslant2$ 的解集是
      
        {\centering\begin{tabular*}{\linewidth}{@{}@{\extracolsep{\fill}}llll@{}}
          \(\text{A. }\{x\mid -2\leqslant x\leqslant1\}\) & \(\text{B. }\{x\mid x\leqslant-2\}\) & \(\text{C. }\{x\mid -2\leqslant x<1\}\) & \(\text{D. }\{x\mid x>1\}\)
        \end{tabular*}\par}

        \item 在 $\triangle ABC$ 中，$BC=2$，$AC=1+\sqrt{3}$，$AB=\sqrt{6}$，则 $\angle A=$
       
        {\centering\begin{tabular*}{\linewidth}{@{}@{\extracolsep{\fill}}llll@{}}
          \(\text{A. }45^\circ\) & \(\text{B. }60^\circ\) & \(\text{C. }120^\circ\) & \(\text{D. }135^\circ\)
        \end{tabular*}\par}

        \item 设抛物线 $C:y^2=2px\,(p>0)$ 的焦点为 $F$，点 $A$ 在 $C$ 上，过 $A$ 作 $C$ 准线的垂线，垂足为 $B$．若直线 $BF$ 的方程为 $y=-2x+2$，则 $|AF|=$
     
        {\centering\begin{tabular*}{\linewidth}{@{}@{\extracolsep{\fill}}llll@{}}
          \(\text{A. }3\) & \(\text{B. }4\) & \(\text{C. }5\) & \(\text{D. }6\)
        \end{tabular*}\par}

        \item 记 $S_n$ 为等差数列 $\{a_n\}$ 的前 $n$ 项和．若 $S_3=6$，$S_5=-5$，则 $S_6=$
      
        {\centering\begin{tabular*}{\linewidth}{@{}@{\extracolsep{\fill}}llll@{}}
          \(\text{A. }-20\) & \(\text{B. }-15\) & \(\text{C. }-10\) & \(\text{D. }-5\)
        \end{tabular*}\par}

        \item 已知 $0<\alpha<\pi$，$\cos\dfrac{\alpha}{2}=\dfrac{\sqrt{5}}{5}$，则 $\sin\bigl(\alpha-\dfrac{\pi}{4}\bigr)=$
      
        {\centering\begin{tabular*}{\linewidth}{@{}@{\extracolsep{\fill}}llll@{}}
          \(\text{A. }\dfrac{\sqrt{2}}{10}\) & \(\text{B. }\dfrac{\sqrt{2}}{5}\) & \(\text{C. }\dfrac{3\sqrt{2}}{10}\) & \(\text{D. }\dfrac{7\sqrt{2}}{10}\)
        \end{tabular*}\par}
    \end{enumerate}
\end{description}

\begin{description}
    \item[二、多选题] 本题共3小题，每小题6分，共18分．在每小题给出的四个选项中，有多项符合题目要求．全部选对的得6分，部分选对的得部分分，有选错的得0分．
    \begin{enumerate}[leftmargin=5pt]
      \addtocounter{enumi}{+8}
        \item 记 $S_n$ 为等比数列 $\{a_n\}$ 的前 $n$ 项和，$q$ 为 $\{a_n\}$ 的公比，$q>0$．若 $S_3=7$，$a_3=1$，则
        \begin{enumerate}[label=\Alph*.]
            \item $q=\dfrac{1}{2}$
            \item $a_5=\dfrac{1}{9}$
            \item $S_5=8$
            \item $a_n+S_n=8$
        \end{enumerate}

        \item 已知 $f(x)$ 是定义在 $\mathbb{R}$ 上的奇函数，且当 $x>0$ 时，$f(x)=(x^2-3)e^x+2$，则
        \begin{enumerate}[label=\Alph*.]
            \item $f(0)=0$
            \item 当 $x<0$ 时，$f(x)=-(x^2-3)e^{-x}-2$
            \item $f(x)\geqslant2$ 当且仅当 $x\geqslant\sqrt{3}$
            \item $x=-1$ 是 $f(x)$ 的极大值点
        \end{enumerate}

        \item 双曲线 $C:\dfrac{x^2}{a^2}-\dfrac{y^2}{b^2}=1\,(a>0,b>0)$ 的左、右焦点分别是 $F_1,F_2$，左、右顶点分别为 $A_1,A_2$，以 $F_1F_2$ 为直径的圆与曲线 $C$ 的一条渐近线交于 $M,N$ 两点，且 $\angle NA_1M=\dfrac{5\pi}{6}$，则
        \begin{enumerate}[label=\Alph*.]
            \item $\angle A_1MA_2=\dfrac{\pi}{6}$
            \item $|MA_1|=2|MA_2|$
            \item $C$ 的离心率为 $\sqrt{13}$
            \item 当 $a=\sqrt{2}$ 时，四边形 $NA_1MA_2$ 的面积为 $8\sqrt{3}$
        \end{enumerate}
    \end{enumerate}
\end{description}

\begin{description}
    \item[三、填空题] 本题共3小题，每小题5分，共15分．
    \begin{enumerate}[leftmargin=5pt]
      \addtocounter{enumi}{+11}
        \item 已知平面向量 $\bm a=(x,1)$，$\bm b=(x-1,2x)$，若 $\bm a\perp(\bm a-\bm b)$，则 $|\bm a|=$ \rule[-0.3ex]{3em}{0.4pt}．
        \item 若 $x=2$ 是函数 $f(x)=(x-1)(x-2)(x-a)$ 的极值点，则 $f(0)=$ \rule[-0.3ex]{3em}{0.4pt}．
        \item 一个底面半径为 $4\,\text{cm}$，高为 $9\,\text{cm}$ 的封闭圆柱形容器（容器壁厚度忽略不计）内有两个半径相等的铁球，则铁球半径的最大值为 \rule[-0.3ex]{3em}{0.4pt} cm．
    \end{enumerate}
\end{description}

\begin{description}
    \item[四、解答题] 本题共5小题，共77分．解答应写出文字说明、证明过程或演算步骤．
    \begin{enumerate}[leftmargin=5pt]
      \addtocounter{enumi}{+14}
        \item （13分）

        已知函数 $f(x)=\cos(2x+\varphi)\,(0\leqslant\varphi<\pi)$，$f(0)=\dfrac{1}{2}$．

        \begin{enumerate}[label=(\arabic*)]
            \item 求 $\varphi$；
            \item 设函数 $g(x)=f(x)+f\bigl(x-\frac{\pi}{6}\bigr)$，求 $g(x)$ 的值域和单调区间．
        \end{enumerate}

        \item （15分）

        已知椭圆 $C:\dfrac{x^2}{a^2}+\dfrac{y^2}{b^2}=1\,(a>b>0)$ 的离心率为 $\dfrac{\sqrt{2}}{2}$，长轴长为 $4$．

        \begin{enumerate}[label=(\arabic*)]
            \item 求 $C$ 的方程；
            \item 过点 $(0,-2)$ 的直线 $l$ 与 $C$ 交于 $A,B$ 两点，$O$ 为坐标原点．若 $\triangle OAB$ 的面积为 $2$，求 $|AB|$．
        \end{enumerate}

        \item （15分）

        如图，在四边形 $ABCD$ 中，$AB\parallel CD$，$\angle DAB=90^\circ$，$F$ 为 $CD$ 的中点，点 $E$ 在 $AB$ 上，$EF\parallel AD$，$AB=3AD$，$CD=2AD$．将四边形 $EFDA$ 沿 $EF$ 翻折至四边形 $EFD'A'$，使得面 $EFD'A'$ 与面 $EFCB$ 所成的二面角为 $60^\circ$．

        \begin{enumerate}[label=(\arabic*)]
            \item 证明：$A'B\parallel$ 平面 $CD'F$；
            \item 求面 $BCD'$ 与面 $EFD'A'$ 所成二面角的正弦值．
        \end{enumerate}

        \item （17分）

        已知函数 $f(x)=\ln(1+x)-x+\dfrac{1}{2}x^2-kx^3$，其中 $0<k<\dfrac{1}{3}$．

        \begin{enumerate}[label=(\arabic*)]
            \item 证明：$f(x)$ 在区间 $(0,+\infty)$ 存在唯一的极值点和唯一的零点；
            \item 设 $x_1,x_2$ 分别为 $f(x)$ 在区间 $(0,+\infty)$ 的极值点和零点．
            \begin{enumerate}[label=(\roman*)]
                \item 设函数 $g(t)=f(x_1+t)-f(x_1-t)$，证明：$g(t)$ 在区间 $(0,x_1)$ 单调递减；
                \item 比较 $2x_1$ 与 $x_2$ 的大小，并证明你的结论．
            \end{enumerate}
        \end{enumerate}

        \item （17分）

        甲、乙两人进行乒乓球练习，每个球胜者得 $1$ 分，负者得 $0$ 分．设每个球甲胜的概率为 $p\,\bigl(\frac{1}{2}<p<1\bigr)$，乙胜的概率为 $q$，$p+q=1$，且各球的胜负相互独立．对正整数 $k\geqslant2$，记 $p_k$ 为打完 $k$ 个球后甲比乙至少多得 $2$ 分的概率，$q_k$ 为打完 $k$ 个球后乙比甲至少多得 $2$ 分的概率．

        \begin{enumerate}[label=(\arabic*)]
            \item 求 $p_3,p_4$（用 $p$ 表示）；
            \item 若 $\dfrac{p_4-p_3}{q_4-q_3}=4$，求 $p$；
            \item 证明：对任意正整数 $m$，$p_{2m+1}-q_{2m+1}<p_{2m}-q_{2m}<p_{2m+2}-q_{2m+2}$．
        \end{enumerate}
    \end{enumerate}
\end{description}

\clearpage


\subsection{2025年北京卷}\label{gk:2025年北京卷}

\begin{description}
    \item[一、选择题] 本题共10小题，每小题4分，共40分．在每小题给出的四个选项中，选出符合题目要求的一项．
    \begin{enumerate}[leftmargin=5pt]
        \item 已知集合 $M=\{x\mid 2x-1>5\}$，$N=\{1,2,3\}$，则 $M\cap N=$
       
        {\centering\begin{tabular*}{\linewidth}{@{}@{\extracolsep{\fill}}llll@{}}
          \(\text{A. }\{1,2,3\}\) & \(\text{B. }\{2,3\}\) & \(\text{C. }\{3\}\) & \(\text{D. }\varnothing\)
        \end{tabular*}\par}

        \item 已知复数 $z$ 满足 $\mathrm{i}\cdot z+2=2\mathrm{i}$，则 $|z|=$
       
        {\centering\begin{tabular*}{\linewidth}{@{}@{\extracolsep{\fill}}llll@{}}
          \(\text{A. }2\) & \(\text{B. }2\sqrt{2}\) & \(\text{C. }4\) & \(\text{D. }8\)
        \end{tabular*}\par}

        \item 双曲线 $x^2-4y^2=4$ 的离心率为
      
        {\centering\begin{tabular*}{\linewidth}{@{}@{\extracolsep{\fill}}llll@{}}
          \(\text{A. }\dfrac{\sqrt{3}}{2}\) & \(\text{B. }\dfrac{\sqrt{5}}{2}\) & \(\text{C. }\dfrac{5}{4}\) & \(\text{D. }\sqrt{5}\)
        \end{tabular*}\par}

        \item 为了得到函数 $y=9^x$ 的图象，只需把函数 $y=3^x$ 的图象上所有点的
      
        {\centering\begin{tabular*}{\linewidth}{@{}@{\extracolsep{\fill}}llll@{}}
          \(\text{A. }\)横坐标变为原来的\(\frac{1}{2}\)倍（纵坐标不变） \\
          \(\text{B. }\)横坐标变为原来的\(2\)倍（纵坐标不变） \\
           \(\text{C. }\)纵坐标变为原来的\(\frac{1}{3}\)倍（横坐标不变） \\
            \(\text{D. }\)纵坐标变为原来的\(3\)倍（横坐标不变）
        \end{tabular*}\par}

        \item 已知 $\{a_n\}$ 是公差不为零的等差数列，$a_1=-2$，若 $a_3,a_4,a_6$ 成等比数列，则 $a_{10}=$
      
        {\centering\begin{tabular*}{\linewidth}{@{}@{\extracolsep{\fill}}llll@{}}
          \(\text{A. }-20\) & \(\text{B. }-18\) & \(\text{C. }16\) & \(\text{D. }18\)
        \end{tabular*}\par}

        \item 已知 $a>0$，$b>0$，则
    
        {\centering\begin{tabular*}{\linewidth}{@{}@{\extracolsep{\fill}}llll@{}}
          \(\text{A. }a^2+b^2>2ab\) & \(\text{B. }\frac{1}{a}+\frac{1}{b}\geqslant\frac{1}{\sqrt{ab}}\) & \(\text{C. }a+b>\sqrt{ab}\) & \(\text{D. }\frac{1}{a}+\frac{1}{b}\leqslant\frac{2}{\sqrt{ab}}\)
        \end{tabular*}\par}

        \item 已知函数 $f(x)$ 的定义域为 $D$，则``$f(x)$ 的值域为 $\mathbb{R}$''是``对任意 $M\in\mathbb{R}$，存在 $x_0\in D$，使得 $|f(x_0)|>M$''的
      
        {\centering\begin{tabular*}{\linewidth}{@{}@{\extracolsep{\fill}}llll@{}}
          \(\text{A. }\)充分不必要条件 & \(\text{B. }\)必要不充分条件 & \(\text{C. }\)充要条件 & \(\text{D. }\)既不充分也不必要条件
        \end{tabular*}\par}

        \item 设函数 $f(x)=\sin\omega x+\cos\omega x\,(\omega>0)$，若 $f(x+\pi)=f(x)$ 恒成立，且 $f(x)$ 在 $\bigl[0,\frac{\pi}{4}\bigr]$ 上存在零点，则 $\omega$ 的最小值为
     
        {\centering\begin{tabular*}{\linewidth}{@{}@{\extracolsep{\fill}}llll@{}}
          \(\text{A. }8\) & \(\text{B. }6\) & \(\text{C. }4\) & \(\text{D. }3\)
        \end{tabular*}\par}

        \item 在一定条件下，某人工智能大语言模型训练 $N$ 个单位的数据量所需要的时间 $T=k\log_2 N$（单位：h），其中 $k$ 为常数．在此条件下，已知训练数据量 $N$ 从 $10^6$ 个单位增加到 $1.024\times10^9$ 个单位时，训练时间增加 $20\,\text{h}$；当训练数据量 $N$ 从 $1.024\times10^9$ 个单位增加到 $4.096\times10^9$ 个单位时，训练时间增加
     
        {\centering\begin{tabular*}{\linewidth}{@{}@{\extracolsep{\fill}}llll@{}}
          \(\text{A. }2\,\text{h}\) & \(\text{B. }4\,\text{h}\) & \(\text{C. }20\,\text{h}\) & \(\text{D. }40\,\text{h}\)
        \end{tabular*}\par}

        \item 在平面直角坐标系 $xOy$ 中，$\bigl|\overrightarrow{OA}\bigr|=\bigl|\overrightarrow{OB}\bigr|=\sqrt{2}$，$\bigl|\overrightarrow{AB}\bigr|=2$，设 $C(3,4)$，则 $\bigl|2\overrightarrow{CA}+\overrightarrow{AB}\bigr|$ 的取值范围是
      
        {\centering\begin{tabular*}{\linewidth}{@{}@{\extracolsep{\fill}}llll@{}}
          \(\text{A. }[6,14]\) & \(\text{B. }[6,12]\) & \(\text{C. }[8,14]\) & \(\text{D. }[8,12]\)
        \end{tabular*}\par}
    \end{enumerate}
\end{description}

\begin{description}
    \item[二、填空题] 本题共5小题，每小题5分，共25分．
    \begin{enumerate}[leftmargin=5pt]
      \addtocounter{enumi}{+10}
        \item 已知抛物线 $y^2=2px\,(p>0)$ 的顶点到焦点的距离为 $3$，则 $p=$ \rule[-0.3ex]{3em}{0.4pt}．
        \item 已知 $(1-2x)^4=a_0-2a_1x+4a_2x^2-8a_3x^3+16a_4x^4$，则 $a_0=$ \uline{\hspace{4em}}；$a_1+a_2+a_3+a_4=$ \rule[-0.3ex]{3em}{0.4pt}．
        \item 已知 $\alpha,\beta\in[0,2\pi]$，且 $\sin(\alpha+\beta)=\sin(\alpha-\beta)$，$\cos(\alpha+\beta)\neq\cos(\alpha-\beta)$．写出满足条件的一组 $\alpha,\beta$ 的值 $\alpha=$ \rule[-0.3ex]{3em}{0.4pt}，$\beta=$ \rule[-0.3ex]{3em}{0.4pt}．
        \item 某科技兴趣小组用3D打印机制作的一个零件可以抽象为如图所示的多面体，其中 $ABCDEF$ 是一个平面多边形，平面 $AFR\perp$ 平面 $ABC$，平面 $CDT\perp$ 平面 $ABC$，$AB\perp BC$，$AB\parallel EF\parallel RS\parallel CD$，$BC\parallel DE\parallel ST\parallel AF$．若 $AB=BC=8$，$AF=CD=4$，$RA=RF=TC=TD=\dfrac{5}{2}$，则该多面体的体积为 \rule[-0.3ex]{3em}{0.4pt}．
        \item 关于定义域为 $\mathbb{R}$ 的函数 $f(x)$，给出下列四个结论：
        \begin{enumerate}[label=\textcircled{\arabic*}]
            \item 存在在 $\mathbb{R}$ 上单调递增的函数 $f(x)$ 使得 $f(x)+f(2x)=-x$ 恒成立；
            \item 存在在 $\mathbb{R}$ 上单调递减的函数 $f(x)$ 使得 $f(x)-f(2x)=x$ 恒成立；
            \item 使得 $f(x)+f(-x)=\cos x$ 恒成立的函数 $f(x)$ 存在且有无穷多个；
            \item 使得 $f(x)-f(-x)=\cos x$ 恒成立的函数 $f(x)$ 存在且有无穷多个．
        \end{enumerate}
        其中正确结论的序号是 \rule[-0.3ex]{3em}{0.4pt}．
    \end{enumerate}
\end{description}

\begin{description}
    \item[三、解答题] 本题共6小题，共85分．解答应写出文字说明、证明过程或演算步骤．
    \begin{enumerate}[leftmargin=5pt]
      \addtocounter{enumi}{+15}
        \item （13分）

        在 $\triangle ABC$ 中，$\cos A=-\dfrac{1}{3}$，$a\sin C=4\sqrt{2}$．

        \begin{enumerate}[label=(\arabic*)]
            \item 求 $c$ 的值；
            \item 再从条件\ding{172}、条件\ding{173}、条件\ding{174}这三个条件中选择一个作为已知，使得 $\triangle ABC$ 存在，求 $BC$ 边上的高．
            \begin{quote}
            条件\ding{172}：$a=6$；条件\ding{173}：$a\sin B=\dfrac{10\sqrt{2}}{3}$；条件\ding{174}：$\triangle ABC$ 的面积为 $10\sqrt{2}$．
            \end{quote}
        \end{enumerate}

        \item （14分）

        如图，在四棱锥 $P-ABCD$ 中，$\triangle ADC$ 与 $\triangle BAC$ 均为等腰直角三角形，$\angle ADC=90^\circ$，$\angle BAC=90^\circ$，$E$ 为 $BC$ 的中点．

        \begin{enumerate}[label=(\arabic*)]
            \item 若 $F,G$ 分别为 $PD,PE$ 的中点，求证：$FG\parallel$ 平面 $PAB$；
            \item 若 $PA\perp$ 平面 $ABCD$，$PA=AC$，求直线 $AB$ 与平面 $PCD$ 所成角的正弦值．
        \end{enumerate}

        \item （14分）

        某次考试中，只有一道单项选择题考查了某个知识点，甲、乙两校的高一年级学生都参加了这次考试．为了解学生对该知识点的掌握情况，随机抽查了甲、乙两校高一年级各 $100$ 名学生该题的答题数据，其中甲校学生选择正确的人数为 $80$，乙校学生选择正确的人数为 $75$．假设学生之间答题相互独立．用频率估计概率．

        \begin{enumerate}[label=(\arabic*)]
            \item 估计甲校高一年级学生该题选择正确的概率 $p$；
            \item 从甲、乙两校高一年级学生中各随机抽取 $1$ 名．设 $X$ 为这 $2$ 名学生中该题选择正确的人数，估计 $X=1$ 的概率及 $X$ 的数学期望；
            \item 假设：如果没有掌握该知识点，学生就从题目给出的四个选项中随机选择一个作为答案；如果掌握该知识点，甲校学生选择正确的概率为 $100\%$，乙校学生选择正确的概率为 $85\%$．设甲、乙两校高一年级学生掌握该知识点的概率估计值分别为 $p_1,p_2$，判断 $p_1$ 与 $p_2$ 的大小．（结论不要求证明）
        \end{enumerate}
        \item （15分）

        已知椭圆 $E:\dfrac{x^2}{a^2}+\dfrac{y^2}{b^2}=1\,(a>b>0)$ 的离心率为 $\dfrac{\sqrt{2}}{2}$，椭圆 $E$ 上的点到两焦点的距离之和为 $4$．

        \begin{enumerate}[label=(\arabic*)]
            \item 求椭圆 $E$ 的方程；
            \item 设 $O$ 为坐标原点，点 $M(x_0,y_0)\,(x_0\neq0)$ 在椭圆 $E$ 上，直线 $x_0x+2y_0y-4=0$ 与直线 $y=2$，$y=-2$ 分别交于点 $A,B$．设 $\triangle OAM$ 与 $\triangle OBM$ 的面积分别为 $S_1,S_2$，比较 $\dfrac{S_1}{S_2}$ 与 $\dfrac{|OA|}{|OB|}$ 的大小．
        \end{enumerate}

        \item （15分）

        已知函数 $f(x)$ 的定义域是 $(-1,+\infty)$，$f(0)=0$，导函数 $f'(x)=\dfrac{\ln(1+x)}{1+x}$．设 $l_1$ 是曲线 $y=f(x)$ 在点 $A(a,f(a))\,(a\neq0)$ 处的切线．

        \begin{enumerate}[label=(\arabic*)]
            \item 求 $f'(x)$ 的最大值；
            \item 当 $-1<a<0$ 时，证明：除切点 $A$ 外，曲线 $y=f(x)$ 在直线 $l_1$ 的上方；
            \item 设过点 $A$ 的直线 $l_2$ 与直线 $l_1$ 垂直，$l_1,l_2$ 与 $x$ 轴交点的横坐标分别是 $x_1,x_2$．若 $a>0$，求 $\dfrac{2a-x_2-x_1}{x_2-x_1}$ 的取值范围．
        \end{enumerate}

        \item （15分）

        已知集合 $A=\{1,2,3,4,5,6,7,8\}$，$M=\{(x,y)\mid x\in A,y\in A\}$．从 $M$ 中选取 $n$ 个不同的元素组成一个序列：$(x_1,y_1),(x_2,y_2),\dots,(x_n,y_n)$，其中 $(x_i,y_i)$ 称为该序列的第 $i$ 项 $(i=1,2,\dots,n)$，若该序列的相邻项 $(x_i,y_i),(x_{i+1},y_{i+1})$ 满足
        \[
        \begin{cases}
        |x_{i+1}-x_i|=3, \\
        |y_{i+1}-y_i|=4
        \end{cases}
        \quad\text{或}\quad
        \begin{cases}
        |x_{i+1}-x_i|=4, \\
        |y_{i+1}-y_i|=3
        \end{cases}
        \quad(i=1,2,\dots,n-1),
        \]
        则称该序列为 $K$ 列．

        \begin{enumerate}[label=(\arabic*)]
            \item 对于第 $1$ 项为 $(3,3)$ 的 $K$ 列，写出它的第 $2$ 项；
            \item 设 $\Gamma$ 为 $K$ 列，且 $\Gamma$ 中的项 $(x_i,y_i)\,(i=1,2,\dots,n)$ 满足：当 $i$ 为奇数时，$x_i\in\{1,2,7,8\}$；当 $i$ 为偶数时，$x_i\in\{3,4,5,6\}$．判断 $(3,2),(4,4)$ 能否同时为 $\Gamma$ 中的项，并说明理由；
            \item 证明：由 $M$ 的全部元素组成的序列都不是 $K$ 列．
        \end{enumerate}
    \end{enumerate}
\end{description}
\clearpage


\subsection{2025年上海卷}\label{gk:2025年上海卷}

\begin{description}
    \item[一、填空题] 本大题共有12题，满分54分，第1\textasciitilde6题每题4分，第7\textasciitilde12题每题5分．
    \begin{enumerate}[leftmargin=5pt]
        \item 已知全集 $U=\{x\mid2\leqslant x\leqslant5,x\in\mathbb{R}\}$，若集合 $A=\{x\mid2\leqslant x<4,x\in\mathbb{R}\}$，则 $\overline{A}=$ \rule[-0.3ex]{3em}{0.4pt}．
        \item 设 $x\in\mathbb{R}$，不等式 $\dfrac{x-1}{x-3}<0$ 的解集为 \rule[-0.3ex]{3em}{0.4pt}．
        \item 若等差数列 $\{a_n\}$ 的首项 $a_1=-3$，公差 $d=2$，则该数列的前 $6$ 项的和为 \rule[-0.3ex]{3em}{0.4pt}．
        \item 在 $(2x-1)^5$ 的二项展开式中，$x^3$ 项的系数为 \rule[-0.3ex]{3em}{0.4pt}．
        \item 函数 $y=\cos x$，$x\in\bigl[-\frac{\pi}{2},\frac{\pi}{4}\bigr]$ 的值域为 \rule[-0.3ex]{3em}{0.4pt}．
        \item 若随机变量 $X$ 的分布为 $\begin{pmatrix}5&6&7\\0.2&0.3&0.5\end{pmatrix}$，则期望 $E(X)=$ \rule[-0.3ex]{3em}{0.4pt}．
        \item 在正四棱柱 $ABCD-A_1B_1C_1D_1$ 中，若 $AC=4\sqrt{2}$，$BD_1=9$，则其体积为 \rule[-0.3ex]{3em}{0.4pt}．
        \item 已知 $a>0,b>0$，则 $\dfrac{a}{b}+\dfrac{b}{a}$ 的最小值为 \rule[-0.3ex]{3em}{0.4pt}．
        \item 有 $4$ 名家长和 $2$ 名儿童去郊游爬山．若 $6$ 个人要排成一队列，要求队列的首和尾均是家长，则不同的排列方式共有 \rule[-0.3ex]{3em}{0.4pt} 种．
        \item $\mathrm{i}$ 为虚数单位，若复数 $z$ 满足 $z^2=(\overline{z})^2$，且 $|z|\leqslant1$，则 $|z-2-3\mathrm{i}|$ 的最小值为 \rule[-0.3ex]{3em}{0.4pt}．
        \item 小申同学观察发现，生活中有时候影子可以完全投射在斜面上．某斜面上有两根长为 $1$ 米且垂直于水平面放置的杆子，与斜面的接触点分别为 $A,B$，它们在阳光的照射下呈现出影子，阳光可视为平行光，其中 $A$ 处杆子的影子在水平面上，长度为 $0.4$ 米，$B$ 处杆子的影子完全在斜面上，长度为 $0.45$ 米．则斜面的坡角 $\theta=$ \rule[-0.3ex]{3em}{0.4pt}．（结果精确到 $0.01^\circ$）
        \item 已知 $f(x)=\begin{cases}0,&x=0,\\-1,&x<0\end{cases}$，向量 $\bm a,\bm b,\bm c$ 是平面内三个不同的单位向量，若 $f(\bm a\cdot\bm b)+f(\bm b\cdot\bm c)+f(\bm c\cdot\bm a)=0$，则 $|\bm a+\bm b+\bm c|$ 的取值范围是 \rule[-0.3ex]{3em}{0.4pt}．
    \end{enumerate}
\end{description}

\begin{description}
    \item[二、选择题] 本题共有4题，满分18分，13、14每题4分，15、16题每题5分．每题有且只有一个正确选项．
    \begin{enumerate}[leftmargin=5pt]
      \addtocounter{enumi}{+12}
        \item 已知事件 $A,B$ 相互独立，若 $P(A)=P(B)=\dfrac{1}{2}$，则 $P(A\cap B)$ 的值为
     
        {\centering\begin{tabular*}{\linewidth}{@{}@{\extracolsep{\fill}}llll@{}}
          \(\text{A. }0\) & \(\text{B. }\dfrac{1}{4}\) & \(\text{C. }\dfrac{1}{2}\) & \(\text{D. }1\)
        \end{tabular*}\par}

        \item 已知 $a>0$ 且 $a\neq1$，$s\in\mathbb{R}$，下列各项中，能推出 $a^s>a$ 的是
    
        {\centering\begin{tabular*}{\linewidth}{@{}@{\extracolsep{\fill}}llll@{}}
          \(\text{A. }a>1,s>0\) & \(\text{B. }a>1,s<0\) & \(\text{C. }0<a<1,s>0\) & \(\text{D. }0<a<1,s<0\)
        \end{tabular*}\par}

        \item 已知点 $A(0,1),B(1,2)$，若动点 $C$ 在曲线 $x^2-y^2=1\,(x\geqslant1,y\geqslant0)$ 上，则 $\triangle ABC$ 的面积
      
        {\centering\begin{tabular*}{\linewidth}{@{}@{\extracolsep{\fill}}llll@{}}
          \(\text{A. }\)有最大值，无最小值 & \(\text{B. }\)无最大值，有最小值 & \(\text{C. }\)有最大值，有最小值 & \(\text{D. }\)无最大值，无最小值
        \end{tabular*}\par}

        \item 设 $\lambda\in\mathbb{R}$，数列 $\{a_n\},\{b_n\},\{c_n\}$ 满足 $a_n=10n-9$，$b_n=2^n$，$c_n=\lambda a_n+(1-\lambda)b_n$．若对任意的 $\lambda\in[0,1]$，$a_n,b_n,c_n$ 均能作为三角形的三边长，则满足条件的正整数 $n$ 有
    
        {\centering\begin{tabular*}{\linewidth}{@{}@{\extracolsep{\fill}}llll@{}}
          \(\text{A. }1\text{个}\) & \(\text{B. }3\text{个}\) & \(\text{C. }4\text{个}\) & \(\text{D. }\)无穷多个
        \end{tabular*}\par}
    \end{enumerate}
\end{description}

\begin{description}
    \item[三、解答题] 本大题共有5题，满分78分．解答下列各题必须写出必要的步骤．
    \begin{enumerate}[leftmargin=5pt]
      \addtocounter{enumi}{+16}
        \item （14分）

        下表为2000年至2024年（偶数年）奥运会男子$4\times100$米混合泳接力项目冠军队的成绩（单位：秒）．
        \begin{table}[H]
        \centering
        \begin{tabular}{|c|c|c|c|c|c|}
        \hline
        年份 & 2000 & 2004 & 2008 & 2012 & 2016 \\
        \hline
        成绩 & 213.73 & 210.68 & 209.34 & 208.92 & 207.22 \\
        \hline
        年份 & 2020 & 2024 & & & \\
        \hline
        成绩 & 206.51 & 205.42 & & & \\
        \hline
        \end{tabular}
        \end{table}

        \begin{enumerate}[label=(\arabic*)]
            \item 求这组数据的极差与中位数；
            \item 从这 $7$ 个数据中任选 $3$ 个，求恰有 $2$ 个数据大于 $211$ 的概率；
            \item 若比赛成绩 $y$ 关于年份 $x$ 的回归方程为 $\hat{y}=-0.311x+\hat{b}$，年份 $x$ 的平均数为 $2006$，请预测 $2028$ 年男子 $4\times100$ 米混合泳接力项目冠军队的成绩．（结果精确到 $0.01$ 秒）
        \end{enumerate}

        \item （14分）

        如图，$P$ 为圆锥的顶点，$O$ 为底面圆心，$AB$ 为底面直径，$\angle APB=\dfrac{\pi}{3}$．

        \begin{enumerate}[label=(\arabic*)]
            \item 若直线 $PA$ 与圆锥底面所成角为 $\dfrac{\pi}{4}$，求圆锥的侧面积；
            \item 若 $Q$ 是母线 $PA$ 的中点，点 $C,D$ 在底面圆周上，$\widehat{AC}$ 弧长为 $\dfrac{\pi}{3}$，$CD\parallel AB$，点 $T$ 在线段 $OC$ 上，求证：直线 $QT\parallel$ 平面 $PBD$．
        \end{enumerate}

        \item （14分）

        设 $m\in\mathbb{R}$，$f(x)=x^2-(m+2)x+m\ln x$．

        \begin{enumerate}[label=(\arabic*)]
            \item 若 $f(1)=0$，求不等式 $f(x)\leqslant x^2-1$ 的解集；
            \item 若 $y=f(x)$ 在区间 $(0,+\infty)$ 上存在极大值，求 $m$ 的取值范围．
        \end{enumerate}

        \item （18分）

        设椭圆 $\Gamma:\dfrac{x^2}{a^2}+\dfrac{y^2}{5}=1\,(a>\sqrt{5})$，点 $A$ 为 $\Gamma$ 的右顶点，点 $M(0,m)\,(m>0)$．

        \begin{enumerate}[label=(\arabic*)]
            \item 若 $\Gamma$ 的一个焦点是 $(2,0)$，求 $\Gamma$ 的离心率；
            \item 设 $a=4$，椭圆 $\Gamma$ 上存在一点 $P$ 满足 $\overrightarrow{PA}=2\overrightarrow{MP}$，求 $m$ 的值；
            \item 若线段 $AM$ 的中垂线 $l$ 的斜率为 $2$，$l$ 与椭圆 $\Gamma$ 交于 $C,D$ 两点，且 $\angle CMD$ 为钝角，求 $a$ 的取值范围．
        \end{enumerate}

        \item （18分）

        已知函数 $y=f(x)$ 的定义域为 $\mathbb{R}$，对于正实数 $a$，定义集合 $M_a=\{x\mid f(x+a)=f(x),x\in\mathbb{R}\}$．

        \begin{enumerate}[label=(\arabic*)]
            \item 设 $f(x)=\sin x$，判断 $\dfrac{\pi}{3}$ 是否是 $M_\pi$ 中的元素，并说明理由；
            \item 设 $f(x)=\begin{cases}x+2,&x<0,\\\sqrt{x},&x\geqslant0,\end{cases}$ 且 $M_a\neq\varnothing$，求 $a$ 的取值范围；
            \item 若 $y=f(x)$ 是偶函数，当 $x\in(0,1]$ 时，$f(x)=1-x$，且对任意 $a\in(0,2)$，均有 $M_a\subseteq M_2$．写出当 $x\in(1,2)$ 时 $f(x)$ 的表达式，并证明：对任意实数 $c$，函数 $y=f(x)-c$ 在区间 $[-3,3]$ 上至多有 $9$ 个零点．
        \end{enumerate}
    \end{enumerate}
\end{description}
\clearpage

\subsection{2025年天津卷}\label{gk:2025年天津卷}

\begin{description}
    \item[一、选择题] 本题共9小题，每小题5分，共45分．在每小题给出的四个选项中，只有一项是符合题目要求的．
    \begin{enumerate}[leftmargin=5pt]
        \item 已知全集 $U=\{1,2,3,4,5\}$，集合 $A=\{1,3\}$，$B=\{2,3,5\}$，则 $\complement_U(A\cup B)=$
       
        {\centering\begin{tabular*}{\linewidth}{@{}@{\extracolsep{\fill}}llll@{}}
          \(\text{A. }\{1,2,4,5\}\) & \(\text{B. }\{1,3,4\}\) & \(\text{C. }\{1,2,3,5\}\) & \(\text{D. }\{4\}\)
        \end{tabular*}\par}

        \item 已知 $x\in\mathbb{R}$，则``$x=0$''是``$\sin2x=0$''的
       
        {\centering\begin{tabular*}{\linewidth}{@{}@{\extracolsep{\fill}}llll@{}}
          \(\text{A. }\)充分不必要条件 & \(\text{B. }\)必要不充分条件 & \(\text{C. }\)充要条件 & \(\text{D. }\)既不充分也不必要条件
        \end{tabular*}\par}

        \item 已知函数 $y=f(x)$ 的部分图象如下，则 $f(x)$ 的解析式可能为
       
        {\centering\begin{tabular*}{\linewidth}{@{}@{\extracolsep{\fill}}llll@{}}
          \(\text{A. }\)\textit{(图象题，选项略)} & \(\text{B. }\)\textit{(图象题)} & \(\text{C. }\)\textit{(图象题)} & \(\text{D. }\)\textit{(图象题)}
        \end{tabular*}\par}

        \item 已知 $m$ 是直线，$\alpha,\beta$ 是两个不同的平面，则下列命题正确的是
     
        {\centering\begin{tabular*}{\linewidth}{@{}@{\extracolsep{\fill}}llll@{}}
          \(\text{A. }\)若 $m\parallel\alpha,m\parallel\beta$，则 $\alpha\parallel\beta$ & \(\text{B. }\)若 $\alpha\perp\beta,m\perp\alpha$，则 $m\parallel\beta$ \\
           \(\text{C. }\)若 $m\subset\alpha,m\perp\beta$，则 $\alpha\perp\beta$ & \(\text{D. }\)若 $\alpha\perp\beta,m\subset\alpha$，则 $m\perp\beta$
        \end{tabular*}\par}

        \item 下列说法错误的是
      
        {\centering\begin{tabular*}{\linewidth}{@{}@{\extracolsep{\fill}}llll@{}}
          \(\text{A. }\)若 $X\sim N(\mu,\sigma^2)$，则 $P(X\leqslant\mu-\sigma)=P(X\geqslant\mu+\sigma)$ \\
           \(\text{B. }\)若 $X\sim N(1,2^2),Y\sim N(2,2^2)$，则 $P(X\leqslant1)<P(Y\leqslant2)$ \\
            \(\text{C. }\)样本相关系数绝对值越接近1，线性相关程度越强 \\
            \(\text{D. }\)样本相关系数绝对值越接近0，线性相关程度越弱
        \end{tabular*}\par}

        \item 已知数列 $\{a_n\}$ 的前 $n$ 项和为 $S_n=-n^2+8n$，则数列 $\{|a_n|\}$ 的前 $12$ 项和为
      
        {\centering\begin{tabular*}{\linewidth}{@{}@{\extracolsep{\fill}}llll@{}}
          \(\text{A. }144\) & \(\text{B. }112\) & \(\text{C. }80\) & \(\text{D. }48\)
        \end{tabular*}\par}

        \item 函数 $f(x)=0.3x-\sqrt{x}$ 的零点所在的一个区间是
      
        {\centering\begin{tabular*}{\linewidth}{@{}@{\extracolsep{\fill}}llll@{}}
          \(\text{A. }(0,0.3)\) & \(\text{B. }(0.3,0.5)\) & \(\text{C. }(0.5,1)\) & \(\text{D. }(1,2)\)
        \end{tabular*}\par}

        \item 已知函数 $f(x)=\sin(\omega x+\varphi)\,(\omega>0,-\pi<\varphi<\pi)$ 在区间 $\bigl[-\frac{5\pi}{12},\frac{\pi}{12}\bigr]$ 上单调递增，直线 $x=\frac{\pi}{12}$ 为曲线 $y=f(x)$ 的一条对称轴，点 $\bigl(\frac{\pi}{12},0\bigr)$ 为曲线 $y=f(x)$ 的一个对称中心，则 $f(x)$ 在区间 $\bigl[0,\frac{\pi}{2}\bigr]$ 上的最小值为
      
        {\centering\begin{tabular*}{\linewidth}{@{}@{\extracolsep{\fill}}llll@{}}
          \(\text{A. }-\dfrac{\sqrt{3}}{2}\) & \(\text{B. }-\dfrac{1}{2}\) & \(\text{C. }-1\) & \(\text{D. }0\)
        \end{tabular*}\par}

        \item 已知双曲线 $\dfrac{x^2}{a^2}-\dfrac{y^2}{b^2}=1\,(a>0,b>0)$ 的左、右焦点分别为 $F_1$ 和 $F_2$，以 $F_2$ 为焦点的抛物线 $y^2=2px\,(p>0)$ 与双曲线在第一象限的交点为 $P$．若 $|PF_1|+|PF_2|=3|F_1F_2|$，则双曲线的离心率为
     
        {\centering\begin{tabular*}{\linewidth}{@{}@{\extracolsep{\fill}}llll@{}}
          \(\text{A. }2\) & \(\text{B. }\sqrt{5}\) & \(\text{C. }\dfrac{\sqrt{2}+1}{2}\) & \(\text{D. }\dfrac{\sqrt{5}+1}{2}\)
        \end{tabular*}\par}
    \end{enumerate}
\end{description}

\begin{description}
    \item[二、填空题] 本题共6小题，每小题5分，共30分．
    \begin{enumerate}[leftmargin=5pt]
      \addtocounter{enumi}{+9}
        \item $\mathrm{i}$ 是虚数单位，$\dfrac{3-\mathrm{i}}{\mathrm{i}}=$ \rule[-0.3ex]{3em}{0.4pt}．
        \item 在 $(x-1)^6$ 的展开式中，$x^3$ 的系数为 \rule[-0.3ex]{3em}{0.4pt}．
        \item 若直线 $x-y+6=0$ 与 $x$ 轴交于点 $A$，与 $y$ 轴交于点 $B$，与圆 $x^2+y^2=r^2$ 相切，则 $r=$ \rule[-0.3ex]{3em}{0.4pt}．
        \item 已知 $a>0,b>0$，且 $a+b=1$，则 $\dfrac{1}{a}+\dfrac{4}{b}$ 的最小值为 \rule[-0.3ex]{3em}{0.4pt}．
        \item 已知 $a,b\in\mathbb{R}$，$\forall x\in[-2,2]$，均有 $(2a+b)x^2+bx-a-1\leqslant0$ 恒成立，求 $2a+b$ 的最小值为 \uline{例\ref{高考:2025年天津卷15}}．
        \item 回忆版暂缺 \rule[-0.3ex]{3em}{0.4pt}．
    \end{enumerate}
\end{description}

\begin{description}
    \item[三、解答题] 本题共5小题，共75分．解答应写出文字说明、证明过程或演算步骤．
    \begin{enumerate}[leftmargin=5pt]
      \addtocounter{enumi}{+15}
        \item （14分）

        在 $\triangle ABC$ 中，角 $A,B,C$ 所对的边分别为 $a,b,c$．已知 $a\sin B=\sqrt{3}b\cos A$，$c-2b=1$，$a=\sqrt{7}$．

        \begin{enumerate}[label=(\arabic*)]
            \item 求角 $A$ 的大小；
            \item 求 $c$ 的值；
            \item 求 $\sin(A+2B)$ 的值．
        \end{enumerate}

        \item （15分）

        如图，在正方体 $ABCD-A_1B_1C_1D_1$ 中，棱长为 $4$，$E,F$ 分别为 $A_1D_1$ 和 $B_1C_1$ 的中点，$G$ 在线段 $CC_1$ 上，且 $CG=3GC_1$．

        \begin{enumerate}[label=(\arabic*)]
            \item 求证：$GF\perp$ 平面 $EBF$；
            \item 求平面 $EBF$ 和平面 $EBG$ 夹角的余弦值；
            \item 求三棱锥 $D-EBF$ 的体积．
        \end{enumerate}

        \item （15分）

        已知椭圆 $\dfrac{x^2}{a^2}+\dfrac{y^2}{b^2}=1\,(a>b>0)$ 的离心率为 $\dfrac{1}{2}$，左焦点为 $F$，右顶点为 $A$，点 $P$ 在直线 $x=a$ 上，直线 $FP$ 的斜率为 $\dfrac{1}{3}$，$\triangle PFA$ 的面积为 $\dfrac{3}{2}$．

        \begin{enumerate}[label=(\arabic*)]
            \item 求椭圆的方程；
            \item 若过 $P$ 的直线与椭圆有唯一公共点 $B$（$B$ 异于 $A$），求证：$FP$ 平分 $\angle AFB$．
        \end{enumerate}

        \item （15分）

        已知等差数列 $\{a_n\}$ 和等比数列 $\{b_n\}$ 满足 $a_1=b_1=2$，$a_2=b_2+1$，$a_3=b_3$．

\begin{enumerate}[label=(\arabic*)]
    \item 求 $\{a_n\}$ 和 $\{b_n\}$ 的通项公式；
    \item 给定 $n\in\mathbb{N}^*$，记集合 $I=\{0,1\}$，集合
    \[
    T_n=\{p_1a_1b_1+p_2a_2b_2+\cdots+p_na_nb_n\mid p_1,p_2,\cdots,p_n\in I\}.
    \]
    \begin{itemize}
        \item[\ding{172}] 求证：对任意的 $t\in T_n$，有 $t<a_{n+1}b_{n+1}$；
        \item[\ding{173}] 求 $T_n$ 中所有元素之和．
    \end{itemize}
\end{enumerate}

        \item （16分）

        已知 $a\in\mathbb{R}$，函数 $f(x)=ax-(\ln x)^2$．

        \begin{enumerate}[label=(\arabic*)]
            \item 当 $a=1$ 时，求曲线 $y=f(x)$ 在点 $(1,f(1))$ 处的切线方程；
            \item 若 $f(x)$ 有三个零点 $x_1,x_2,x_3$，且 $x_1<x_2<x_3$，
            \begin{enumerate}[label=\textcircled{\arabic*}]
                \item 求 $a$ 的取值范围；
                \item 求证：$(\ln x_2-\ln x_1)\ln x_3<\dfrac{4e}{e-1}$．
            \end{enumerate}
        \end{enumerate}
    \end{enumerate}
\end{description}

